\chapter{Introduction}
\label{ch:introduction}

Personal journaling has evolved from simple note-taking into a practice supported by digital tools that encourage self-reflection and emotional awareness. Recent advances in mobile hardware and on-device machine learning now allow private analysis of personal text without cloud services. This thesis examines whether those capabilities can support regular journalers who want a clearer picture of their emotional patterns over time.

The project proposes a mobile microjournaling system that processes short user entries on the device using compact language models, emotion detection, and clustering. Semantic representations of journal entries are grouped by similarity to surface recurring themes and emotional trends not visible in a flat chronological list. All computation runs locally, so user text never leaves the device.

The methodology builds on prior work in emotion recognition on edge devices \cite{Pandelea2021EmotionEdge}, efficient clustering for mobile processors \cite{awad2022improved}, and NLP-assisted journaling systems \cite{10351450}. Compact transformer architectures such as MobileBERT \cite{sun2020mobilebert}, DistilBERT \cite{sanh2019distilbert}, and TinyBERT \cite{jiao2020tinybert} show that competitive language understanding is achievable at a fraction of full-size BERT's cost, well within the budget of a smartphone NPU.

A more recent development relevant to on-device deployment is Matryoshka Representation Learning (MRL) \cite{kusupati2022matryoshka}. MRL trains a model with a nested loss so that every prefix sub-vector of the output is itself a useful embedding. In practice this means a single model can produce vectors at multiple dimensionalities without retraining or recompression: a 768-dimensional output can be truncated to 256 or 128 dimensions at inference time. This thesis evaluates nomic-embed-text-v1.5 \cite{nussbaum2024nomic}, an MRL-trained model, alongside all-MiniLM-L6-v2 \cite{reimers2019sentencebert} and Apple's NLEmbedding, to understand whether MRL dimension flexibility translates into a practical trade-off for on-device journaling.

The central question this thesis investigates is: at what MRL truncation level does Nomic match all-MiniLM-L6-v2 on journal-entry clustering quality, and what does the dimension-versus-quality curve look like across standard and domain-specific benchmarks? The short answer, established in Chapter~\ref{chap:evaluation}, is that 256-dimensional Nomic embeddings retain nearly all of the Spearman correlation relative to the full 768-dimensional model, but MiniLM still leads on clustering V-measure when tested against actual journal entries.

The thesis is organised as follows. Chapter~\ref{ch:background} reviews related work on digital journaling, on-device NLP, sentence embedding models, MRL, and mobile hardware. Chapter~\ref{ch:requirements} specifies the functional and non-functional requirements for GJournal and defines the architecture. Chapter~\ref{ch:implementation} describes the implementation: the data model, AI pipeline, clustering algorithm, and user interface. Chapter~\ref{chap:evaluation} presents the empirical evaluation across four models on STS-B, TwentyNewsgroups clustering, domain validation, and an on-device efficiency benchmark. Chapter~\ref{conclusions} summarises the findings and outlines directions for future work.

\section*{Note on the use of generative AI}
During the preparation of this thesis, the author used Claude Opus to improve the language and readability of the written text and to improve the readability of the source code. After using Claude Opus, the author reviewed and edited the output as needed. All experimental results, model training, system design, evaluation, and scientific conclusions are entirely the author's own work. The author takes full responsibility for the accuracy and integrity of the final document and codebase.
